Optical lattice and single-ion clocks now reach fractional systematic
uncertainties at or below $1\times10^{-18}$~\citep{arnold2026lu,
zhang2026ca}, and two-clock instabilities approach $10^{-18}$ within an
hour~\citep{oelker2019stability}. Their value compounds when they are
compared. Comparing independent clocks verifies accuracy budgets that no
single laboratory can test, because the environmental and geodetic
assumptions entering a self-evaluation are shared within one site; linking
clocks at different locations into a cooperative \emph{network} is
therefore how optical clocks become verifiable metrological
infrastructure~\citep{riehle2017networks, komar2014quantumnetwork}. The
idea has moved from proposal to practice: fibre links with added noise below
the clocks themselves now span thousands of
kilometres~\citep{schioppo2022cavities, chen2026dissemination}, and
multi-clock campaigns have measured tens of ratios among
clocks in six countries~\citep{lindvall2025coordinated} and across a
continental network~\citep{pizzocaro2026european}.

The most immediate purpose of such a network is the redefinition of the SI
second. The roadmap towards the redefinition requires, among its mandatory
criteria, that frequency ratios be measured independently in different
laboratories with an agreement better than \refThreshold, and that
individual ratios be repeated by different
institutes~\citep{dimarcq2024roadmap}; the requirement was restated in the
2025 CCTF progress
reports~\citep{cctf2025strategy, cctf2025taskforce}. The present status is
asymmetric. Ratios at or below the threshold have been demonstrated only
between clocks on the same campus, most recently over a
\refSameCampusKm~km stabilised link with uncertainties at or below
\refSameCampusU~\citep{aeppli2026nist}; comparisons across independent
institutions have reached \refEuBestU at
best~\citep{pizzocaro2026european}. No frequency-ratio measurement between
clocks in different cities has yet met the
requirement~\citep{fortier2026review}.

Closing that gap matters beyond the definition itself. Chronometrically
levelled clocks map the gravitational potential of the Earth at centimetre
resolution~\citep{grotti2018geodesy, takamoto2020redshift, mcgrew2018geodesy};
clock networks have been proposed and used as sensors for dark
matter~\citep{derevianko2014topological, wcislo2018global}; and
intercontinental clock comparisons through very long baseline
interferometry (VLBI) have linked clocks $\refVlbiKm$\,km apart in Italy
and Japan~\citep{pizzocaro2021vlbi}. All of these applications assume that
remote clocks agree at a level that, until now, has only been demonstrated
locally.

Here we report an inter-city, cross-species frequency-ratio measurement
between an IAPMST ytterbium lattice clock in Wuhan and a USTC strontium
lattice clock in Shanghai, about $\siteKm$\,km apart, connected by a
$\fibreKm$\,km phase-stabilised fibre link and operated coherently for
\nDays~days in three stages. From \nSeg~retained data segments
(\nSamples~seconds), we determine the Yb/Sr ratio to a total standard
uncertainty of \uTotalShort$\times10^{-18}$, meeting the uncertainty
requirement of the redefinition roadmap; the time-varying
gravitational-redshift term is carried as an uncertainty item rather than
applied as a correction. The result establishes the two sites as a
prototype node pair for a distributed optical-clock network, and completes,
at inter-city scale, the class of ratio measurements that the redefinition
process depends on.
